% ============================================================
%  CAPÍTULO 5 — DISEÑO DE LA SOLUCIÓN
%  TFG: GlucoMate — Aplicación de Monitorización Continua de Glucosa
%  Ingeniería de Software
% ============================================================

\chapter{Diseño de la Solución}
\label{cap:diseno}

Este capítulo describe las decisiones de diseño que dan forma al sistema GlucoMate:
la arquitectura general, la descomposición en módulos, el esquema de base de datos,
el diseño de la API REST, la visualización de datos y el diseño del agente de
inteligencia artificial.

% ──────────────────────────────────────────────────────────────
\section{Arquitectura general del sistema}
\label{sec:arquitectura-general}

GlucoMate sigue una \textbf{arquitectura cliente-servidor en tres capas} que separa
de forma limpia la presentación, la lógica de negocio y el almacenamiento de datos.
La figura~\ref{fig:arquitectura-general} muestra el diagrama completo del sistema
con todos sus componentes y los canales de comunicación entre ellos.

\begin{figure}[H]
\centering
\resizebox{0.88\textwidth}{!}{%
\begin{tikzpicture}[
    box/.style={draw, rounded corners, align=center, minimum width=2.9cm, minimum height=0.9cm, font=\small},
    ext/.style={box, fill=orange!12, draw=orange!70!black},
    app/.style={box, fill=blue!10, draw=blue!60!black},
    srv/.style={box, fill=green!10, draw=green!50!black},
    db/.style={box, fill=gray!15, draw=black!60},
    label/.style={font=\scriptsize, fill=white, inner sep=1pt}
]
    \node[ext] (sensor) at (0,0) {Sensor\\Freestyle Libre 2 Plus};
    \node[ext] (juggluco) at (0,-1.8) {Juggluco\\puente local};
    \node[app] (frontend) at (0,-3.6) {Frontend móvil\\React Native + Expo};
    \node[srv] (backend) at (-3.6,-5.6) {Backend REST\\Node.js + Express};
    \node[db] (postgres) at (3.6,-5.6) {PostgreSQL\\persistencia};
    \node[ext] (llama) at (0,-7.3) {\texttt{llama-3.3-70b-versatile}\\Groq API};
    \draw[-{Latex[length=2.6mm]}, thick] (sensor) -- node[label, right] {BLE / NFC} (juggluco);
    \draw[-{Latex[length=2.6mm]}, thick] (juggluco) -- node[label, right] {HTTP local} (frontend);
    \draw[-{Latex[length=2.6mm]}, thick] (frontend) -- node[label, below left] {API REST} (backend);
    \draw[-{Latex[length=2.6mm]}, thick] (frontend) -- node[label, below right] {histórico} (postgres);
    \draw[-{Latex[length=2.6mm]}, thick] (backend) -- node[label, below] {SQL} (postgres);
    \draw[-{Latex[length=2.6mm]}, thick] (backend) -- node[label, left] {prompt} (llama);
\end{tikzpicture}%
}

\caption{Arquitectura general del sistema GlucoMate. El sensor se comunica con
         Juggluco por BLE; el frontend obtiene los datos vía HTTP local y los
         persiste a través de la API REST; el backend accede a PostgreSQL y a
         la API de Groq para consumir el modelo \texttt{llama-3.3-70b-versatile}.}
\label{fig:arquitectura-general}
\end{figure}

\noindent La arquitectura presenta las siguientes características destacables:

\begin{itemize}
    \item \textbf{Desacoplamiento completo entre capas.} El frontend no accede
          directamente a la base de datos; toda la lógica de datos pasa por la
          API REST del backend. Esto permite actualizar el esquema de base de
          datos sin modificar el cliente.

    \item \textbf{Juggluco como puente de seguridad.} El sensor Freestyle Libre 2
          Plus utiliza un enlace BLE cifrado con claves derivadas del número de
          serie del sensor. Juggluco gestiona ese descifrado y expone los datos
          ya procesados en mg/dL a través de un servidor HTTP en \texttt{localhost},
          eliminando la necesidad de reimplementar la criptografía propietaria de
          Abbott en GlucoMate.

    \item \textbf{IA en el backend.} La integración con Groq y el modelo
          \texttt{llama-3.3-70b-versatile} se realiza
          exclusivamente en el servidor, garantizando que la clave de API nunca
          se distribuye con el APK de la aplicación.

    \item \textbf{Operación en red local.} En la versión actual, el frontend y el
          backend operan en la misma red Wi-Fi local, con la IP del servidor
          configurada estáticamente. Esto es suficiente para la validación del
          prototipo y establece las bases para una futura migración a infraestructura
          en la nube.
\end{itemize}

% ──────────────────────────────────────────────────────────────
\section{Diseño de la aplicación móvil (React Native + Expo)}
\label{sec:arquitectura-frontend}

El frontend de GlucoMate está estructurado como una \textbf{Single Screen Application
con navegación por pestañas}, implementada con \texttt{@react-navigation/bottom-tabs}.
La figura~\ref{fig:arquitectura-frontend} muestra la jerarquía de componentes.

\begin{figure}[H]
\centering
\resizebox{0.92\textwidth}{!}{%
\begin{tikzpicture}[
    comp/.style={draw, rounded corners, align=center, minimum width=2.6cm, minimum height=0.8cm, font=\small},
    root/.style={comp, fill=blue!10, draw=blue!60!black},
    screen/.style={comp, fill=green!10, draw=green!50!black},
    child/.style={comp, fill=gray!10, draw=black!55, minimum width=2.4cm, minimum height=0.62cm, font=\scriptsize}
]
    \node[root] (app) at (0,0) {App.js\\NavigationContainer};
    \node[screen] (glucose) at (-4.0,-1.9) {Pestaña Glucosa\\\texttt{NFCScanner.js}};
    \node[screen] (profile) at (4.0,-1.9) {Pestaña Mi Perfil\\\texttt{ProfileScreen.js}};
    \node[child] (graph) at (-8.4,-3.7) {renderGraph + D3};
    \node[child] (food) at (-4.0,-3.7) {FoodModal};
    \node[child] (exercise) at (-0.6,-3.7) {ExerciseModal};
    \node[child] (chat) at (-6.4,-5.5) {ChatModal};
    \node[child] (insulin) at (-2.0,-5.5) {InsulinModal};
    \node[child] (picker) at (2.0,-4.7) {Selector de insulinas};
    \node[child] (icr) at (4.0,-3.7) {Ratios ICR};
    \node[child] (fsi) at (6.0,-4.7) {Factor FSI};
    \draw[-{Latex[length=2.4mm]}, thick] (app.south west) -- (glucose.north);
    \draw[-{Latex[length=2.4mm]}, thick] (app.south east) -- (profile.north);
    \draw[-{Latex[length=2.2mm]}, thick] (glucose.south west) -- (graph.north);
    \draw[-{Latex[length=2.2mm]}, thick] (glucose.south) -- (food.north);
    \draw[-{Latex[length=2.2mm]}, thick] (glucose.south east) -- (exercise.north);
    \draw[-{Latex[length=2.2mm]}, thick] (glucose.south east) -- (insulin.north);
    \draw[-{Latex[length=2.2mm]}, thick] (glucose.south west) -- (chat.north);
    \draw[-{Latex[length=2.2mm]}, thick] (profile.south west) -- (picker.north);
    \draw[-{Latex[length=2.2mm]}, thick] (profile.south) -- (icr.north);
    \draw[-{Latex[length=2.2mm]}, thick] (profile.south east) -- (fsi.north);
\end{tikzpicture}%
}
\caption{Jerarquía de componentes del frontend de GlucoMate.}
\label{fig:arquitectura-frontend}
\end{figure}

\texttt{NFCScanner.js} actúa como el \textbf{componente raíz de la pestaña principal}
y concentra la mayor parte del estado de la aplicación mediante los hooks de React
\texttt{useState} y \texttt{useEffect}. Es el responsable de:

\begin{itemize}
    \item Inicializar el gestor NFC al montar el componente y cancelar la solicitud
          de tecnología al desmontar, evitando fugas de recursos.
    \item Cargar el historial de glucosa, el perfil de usuario y los eventos clínicos
          de las últimas 12 horas en paralelo.
    \item Orquestar el flujo de escaneo NFC → llamada a Juggluco → persistencia en
          backend → recarga de historial.
    \item Gestionar la visibilidad de los cuatro modales (insulina, comida, ejercicio,
          chat IA) mediante flags.
\end{itemize}

% ──────────────────────────────────────────────────────────────
\section{Diseño de la base de datos}
\label{sec:diseno-bd}

El esquema de la base de datos PostgreSQL está diseñado para modelar los conceptos
clínicos relevantes en el autocontrol de la diabetes tipo 1: lecturas de glucosa,
perfil de tratamiento individual, eventos de insulina, comidas y ejercicio.

\subsection{Diagrama Entidad-Relación}
\label{subsec:er-diagram}

\begin{figure}[H]
\centering
\resizebox{0.92\textwidth}{!}{%
\begin{tikzpicture}[
    entity/.style={draw, rounded corners, align=left, minimum width=2.7cm, text width=2.5cm, font=\scriptsize, inner sep=4pt},
    rel/.style={-{Latex[length=2.2mm]}, thick}
]
    \node[entity, fill=blue!8] (glucose) at (-4.2,0) {\textbf{glucose\_readings}\\id (PK)\\user\_id\\glucose\_value\\trend\\reading\_time\\created\_at};
    \node[entity, fill=green!8] (profile) at (0,0) {\textbf{user\_profile}\\id (PK)\\fast\_insulin\_id\\slow\_insulin\_id\\ICR + FSI\\updated\_at};
    \node[entity, fill=yellow!10] (insulins) at (4.2,0) {\textbf{insulins}\\id (PK)\\name\\type\\duration\_hours\\peak\_hours};
    \node[entity, fill=orange!10] (food) at (-4.2,-3.1) {\textbf{food\_events}\\id (PK)\\user\_id\\carbs\_g\\event\_time\\created\_at};
    \node[entity, fill=orange!10] (insEvents) at (0,-3.1) {\textbf{insulin\_events}\\id (PK)\\user\_id\\insulin\_id\\units\\event\_time\\created\_at};
    \node[entity, fill=orange!10] (exercise) at (4.2,-3.1) {\textbf{exercise\_events}\\id (PK)\\user\_id\\sport\_id\\start\_time\\duration\_minutes\\created\_at};
    \node[entity, fill=purple!10] (sports) at (4.2,-6.0) {\textbf{sports}\\id (PK)\\name\\intensity\_type\\danger\_window};
    \draw[rel] (profile) -- node[above, font=\scriptsize] {configura} (insulins);
    \draw[rel] (insEvents) -- node[right, font=\scriptsize] {usa} (insulins);
    \draw[rel] (exercise) -- node[right, font=\scriptsize] {usa} (sports);
    \draw[rel] (glucose) -- node[above, font=\scriptsize] {contexto} (profile);
    \draw[rel] (food) -- node[above, font=\scriptsize] {registro} (insEvents);
\end{tikzpicture}%
}
\caption{Diagrama entidad-relación simplificado del esquema de base de datos de
         GlucoMate.}
\label{fig:er-diagram}
\end{figure}

\subsection{Descripción de las tablas}
\label{subsec:tablas}

\begin{description}
    \item[\texttt{glucose\_readings}] Almacena cada lectura de glucosa recibida del
    sensor. El campo \texttt{trend} recoge la dirección de la tendencia tal como la
    reporta Juggluco (\texttt{Flat}, \texttt{SingleUp}, \texttt{FortyFiveDown}, etc.),
    y \texttt{reading\_time} es un \texttt{TIMESTAMP} que permite las consultas de
    ventana temporal (\texttt{NOW() - INTERVAL '12 HOURS'}) usadas por la gráfica.

    \item[\texttt{insulins}] Catálogo de insulinas conocidas, con su tipo
    (\texttt{fast}/\texttt{slow}), duración total en horas (\texttt{duration\_hours}) y
    hora de pico de acción (\texttt{peak\_hours}). Esta tabla es referenciada por
    \texttt{user\_profile} y por \texttt{insulin\_events}, y sus valores se utilizan
    en la gráfica para calcular las ventanas de insulina activa.

    \item[\texttt{user\_profile}] Contiene exactamente una fila (con \texttt{id = 1})
    que representa el perfil de tratamiento del usuario. Los campos \texttt{icr\_*}
    son los ratios insulina-carbohidrato para cada franja horaria del día
    (desayuno, media mañana, almuerzo, merienda y cena), y \texttt{fsi} es el
    Factor de Sensibilidad a la Insulina (mg/dL por unidad). La tabla se gestiona
    mediante un \texttt{UPSERT} (\texttt{INSERT ... ON CONFLICT DO UPDATE}) que
    garantiza que siempre existe una única fila, independientemente del número de
    veces que se guarde el perfil.

    \item[\texttt{insulin\_events}] Registra cada inyección de insulina realizada por
    el usuario. La clave foránea \texttt{insulin\_id} permite conocer la duración y el
    pico de cada dosis, datos necesarios para calcular la insulina activa (IOB) y para
    renderizar las ventanas rosas en la gráfica.

    \item[\texttt{food\_events}] Registra cada comida con sus gramos de carbohidratos
    y la hora de ingesta. Esta información aparece en la gráfica como marcadores
    verdes con la cantidad de hidratos, y es el dato de entrada principal para el
    cálculo de dosis del agente de IA.

    \item[\texttt{sports}] Catálogo de actividades deportivas con su tipo de
    intensidad (\texttt{Aeróbico}, \texttt{Mixto}, \texttt{Baja intensidad},
    \texttt{Anaeróbico}) y su ventana de peligro en horas
    (\texttt{danger\_window\_hours}), que determina cuánto tiempo después del
    ejercicio puede seguir afectando a la glucosa.

    \item[\texttt{exercise\_events}] Registra cada sesión de ejercicio con su
    deporte, hora de inicio y duración en minutos. Esta información se usa tanto
    para la representación gráfica (ventana naranja con línea de peligro
    postejercicio punteada) como para el contexto del agente de IA, que aplica
    reducciones de bolo en función del tipo de intensidad del deporte.
\end{description}

\subsection{Decisiones de diseño relevantes de la BD}
\label{subsec:decisiones-bd}

\begin{itemize}
    \item \textbf{Perfil de usuario como fila única.} El diseño de \texttt{user\_profile}
          con \texttt{id = 1} fijo es una simplificación deliberada para un sistema
          monousuario. El mecanismo \texttt{UPSERT} garantiza que la pantalla de perfil
          siempre opera correctamente, tanto en la primera configuración como en
          actualizaciones posteriores.

    \item \textbf{Catálogos de insulinas y deportes.} Separar las propiedades
          farmacológicas de las insulinas (duración, pico) y las características de los
          deportes (intensidad, ventana de peligro) en tablas de catálogo independientes
          permite que el sistema sea extensible: añadir una nueva insulina o un nuevo
          deporte no requiere cambios en la lógica de la aplicación.

    \item \textbf{Transacción atómica en registro de comida.} Cuando el usuario
          registra una comida e indica también un bolo de insulina, el backend ejecuta
          ambas inserciones (\texttt{food\_events} e \texttt{insulin\_events}) dentro
          de una misma transacción. Si cualquiera de las dos falla, se hace
          \texttt{ROLLBACK} completo, garantizando que nunca quede una comida registrada
          sin su insulina asociada ni viceversa.
\end{itemize}

% ──────────────────────────────────────────────────────────────
\section{Diseño de la API REST}
\label{sec:diseno-api}

El backend expone una API REST que constituye el único canal de comunicación entre
el frontend y los datos. La tabla~\ref{tab:api-endpoints} recoge todos los
\textit{endpoints} implementados.

\begin{table}[H]
\centering
\caption{Endpoints de la API REST de GlucoMate.}
\label{tab:api-endpoints}
\begin{tabular}{|l|l|p{7cm}|}
\hline
\textbf{Método} & \textbf{Ruta} & \textbf{Descripción} \\
\hline
\texttt{POST} & \texttt{/api/glucose}            & Persiste una nueva lectura de glucosa (valor, tendencia, timestamp). \\
\texttt{GET}  & \texttt{/api/glucose/history}    & Devuelve las lecturas de las últimas 12 horas ordenadas por tiempo, para la gráfica D3. \\
\hline
\texttt{GET}  & \texttt{/api/insulins}           & Devuelve el catálogo completo de insulinas (rápidas y lentas). \\
\texttt{GET}  & \texttt{/api/profile}            & Devuelve el perfil de tratamiento del usuario (JOIN con insulinas). \\
\texttt{PUT}  & \texttt{/api/profile}            & Crea o actualiza el perfil de tratamiento (UPSERT). \\
\hline
\texttt{GET}  & \texttt{/api/insulin-events}     & Devuelve inyecciones de insulina recientes (parámetro \texttt{hours}, por defecto 24). \\
\texttt{POST} & \texttt{/api/insulin-events}     & Registra una nueva inyección de insulina. \\
\hline
\texttt{GET}  & \texttt{/api/food-events}        & Devuelve comidas registradas recientes (parámetro \texttt{hours}). \\
\texttt{POST} & \texttt{/api/food-events}        & Registra una comida y, opcionalmente, el bolo de insulina asociado (transacción). \\
\hline
\texttt{GET}  & \texttt{/api/sports}             & Devuelve el catálogo de deportes disponibles. \\
\texttt{GET}  & \texttt{/api/exercise-events}    & Devuelve sesiones de ejercicio recientes. \\
\texttt{POST} & \texttt{/api/exercise-events}    & Registra una nueva sesión de ejercicio. \\
\hline
\texttt{POST} & \texttt{/api/chat}               & Envía un mensaje al agente de IA. Recopila contexto en tiempo real, construye el \textit{system prompt} y llama a Groq para inferir con \texttt{llama-3.3-70b-versatile}. \\
\hline
\end{tabular}
\end{table}

Todos los \textit{endpoints} de escritura aceptan y devuelven JSON con cabecera
\texttt{Content-Type: application/json}. Las respuestas de error incluyen siempre un
campo \texttt{error} con un mensaje descriptivo y un código HTTP apropiado (400 para
errores de validación, 500 para errores internos del servidor).

% ──────────────────────────────────────────────────────────────
\section{Decisiones de diseño}
\label{sec:diseno-grafica}

La gráfica de glucosa es el elemento central de la interfaz de GlucoMate y el que
mayor complejidad de diseño presenta. Se han tomado las siguientes decisiones de diseño:

\subsection{Ventana temporal y escalas}
\label{subsec:ventana-temporal}

La gráfica muestra siempre una ventana de \textbf{12 horas}: 11 horas hacia el pasado
desde el momento actual y 1 hora hacia el futuro. Esta decisión reproduce la convención
clínica establecida por las principales aplicaciones de MCG (Dexcom CLARITY, LibreLinkUp)
y permite al usuario ver la tendencia reciente con suficiente contexto histórico.

Las escalas se calculan dinámicamente en cada renderizado:

\begin{itemize}
    \item \textbf{Eje X (temporal):} \texttt{d3Scale.scaleTime()} mapea el rango de
          timestamps [\texttt{ahora - 11h}, \texttt{ahora + 1h}] al ancho de pantalla
          en píxeles, con un margen lateral de 20px en cada extremo.

    \item \textbf{Eje Y (glucosa):} \texttt{d3Scale.scaleLinear()} mapea el rango
          [50, \texttt{max\_12h + 50}] a la altura del SVG, con el valor mínimo fijo
          en 50 mg/dL y el máximo calculado dinámicamente como el pico del historial
          más un margen de 50 mg/dL para evitar que la curva toque el borde superior.
\end{itemize}

\subsection{Curva de glucosa}
\label{subsec:curva}

La curva principal se genera con \texttt{d3Shape.line().curve(d3Shape.curveMonotoneX)},
que produce una interpolación cúbica monotónica. Esta elección sobre otras curvas de D3
(cardinal, catmull-rom, natural) se justifica porque \texttt{curveMonotoneX} garantiza
que la curva no oscila artificialmente entre puntos: si los datos van hacia arriba, la
curva va hacia arriba, lo que es crítico en un contexto clínico donde una oscilación
falsa podría malinterpretarse como un evento real de hipoglucemia o hiperglucemia.

\subsection{Líneas de referencia clínicas}
\label{subsec:lineas-referencia}

La gráfica incluye dos líneas horizontales de referencia punteadas en los valores
\textbf{70 mg/dL} (umbral de hipoglucemia) y \textbf{180 mg/dL} (umbral de
hiperglucemia), con la zona entre ellas coloreada en verde translúcido
(\texttt{rgba(200, 230, 201, 0.4)}) como indicador visual del rango objetivo.
Estos valores son los estándares clínicos reconocidos por la American Diabetes
Association (ADA).

\subsection{Capas de eventos superpuestas}
\label{subsec:capas-eventos}

Una de las características más innovadoras de la gráfica de GlucoMate es la
superposición de múltiples capas de información clínica sobre la curva de glucosa.
El diseño contempla tres capas independientes, cada una activable/desactivable
mediante los \textit{chips} de filtro:

\begin{description}
    \item[\textbf{Capa de insulina rápida (rosa).}] Para cada inyección de insulina
    rápida activa, se dibuja un rectángulo translúcido (\texttt{rgba(255, 105, 180, 0.3)})
    que abarca desde el momento de la inyección hasta el fin de su duración. Una línea
    vertical punteada de color rosa intenso (\texttt{\#e91e63}) marca el momento de
    pico de acción de la insulina, el instante en que su efecto hipoglucemiante es
    máximo.

    \item[\textbf{Capa de ejercicio (naranja).}] Para cada sesión de ejercicio, se
    dibuja un rectángulo translúcido (\texttt{rgba(255, 152, 0, 0.3)}) durante la
    duración de la actividad. Una línea punteada naranja horizontal indica la
    extensión de la ventana de peligro postejercicio (\texttt{danger\_window\_hours}),
    el periodo durante el cual el ejercicio sigue afectando a la sensibilidad a la
    insulina.

    \item[\textbf{Capa de comidas (verde).}] Cada comida se representa con un
    marcador circular verde en la parte inferior de la gráfica, mostrando los gramos
    de carbohidratos registrados, y una línea vertical punteada que sube hasta la
    parte superior del área de dibujo.


\end{description}

Una capa superpuesta sobre la gráfica pero siempre visible y no desactivable es las\textbf{etiquetas de eje X.} Los ticks del eje temporal se muestran cada
    dos horas en horas pares (00:00, 02:00, 04:00...), calculadas de forma dinámica
    para que siempre sean horas redondas independientemente del momento del día.

\subsection{Indicador de glucosa actual}
\label{subsec:indicador-glucosa}

El valor actual de glucosa se muestra con una tipografía de \textbf{110 puntos} en la
parte superior de la pantalla, coloreado dinámicamente según el rango:

\begin{itemize}
    \item \textcolor{blue}{\textbf{Azul}} (\texttt{\#2196F3}): glucosa en rango normal (70--180 mg/dL).
    \item \textcolor{red}{\textbf{Rojo}} (\texttt{\#f44336}): hipoglucemia (menor a 70 mg/dL).
    \item \textcolor{orange}{\textbf{Naranja}} (\texttt{\#ff9800}): hiperglucemia (mayor a 180 mg/dL).
    \item \textbf{Gris} (\texttt{\#9e9e9e}): dato antiguo o sin lectura reciente.
\end{itemize}

La flecha de tendencia (\texttt{up / up-right / right / down-right / down}) se renderiza con una fuente
personalizada (\texttt{ArrowFont}) cargada mediante \texttt{expo-font}, ya que los
caracteres de flecha estándar de Unicode no tienen representación tipográfica uniforme
en todos los dispositivos Android.

% ──────────────────────────────────────────────────────────────
\section{Diseño de la capa de presentación (UI/UX)}
\label{sec:diseno-modales}

Los tres modales de registro de eventos clínicos (insulina, comida y ejercicio) siguen
un patrón de diseño uniforme:

\begin{enumerate}
    \item \textbf{Modal superpuesto con overlay semitransparente.} Cada modal se
          presenta como un \texttt{Modal} de React Native con \texttt{transparent=true}
          y un overlay \texttt{rgba(0,0,0,0.5)} que oscurece el contenido subyacente,
          manteniendo el contexto visual de la pantalla principal.

    \item \textbf{Selector de hora manual.} Los tres modales incluyen un selector
          de hora compuesto por dos campos \texttt{TextInput} (horas y minutos),
          inicializados con la hora actual del dispositivo. Si la hora introducida
          es posterior a la hora actual del sistema, el modal la interpreta
          automáticamente como perteneciente al día anterior, resolviendo el caso
          de uso habitual de registrar una inyección nocturna antes de medianoche.

    \item \textbf{Feedback visual de estado.} El botón de guardar muestra un
          \texttt{ActivityIndicator} durante el guardado y queda deshabilitado
          hasta que los campos obligatorios son válidos, previniendo envíos duplicados.
\end{enumerate}

Las particularidades de cada modal son:

\begin{description}
    \item[\texttt{InsulinModal.js}] Permite seleccionar entre insulina rápida e
    insulina lenta mediante un selector de dos botones cuyo estado activo muestra
    los colores específicos de cada tipo (rosa para rápida, morado para lenta),
    tomando los nombres de insulina configurados en el perfil del usuario.

    \item[\texttt{FoodModal.js}] Incluye dos campos numéricos: gramos de
    carbohidratos y unidades de bolo de insulina rápida. El campo de bolo queda
    deshabilitado si el usuario no tiene insulina rápida configurada en su perfil,
    mostrando un mensaje de error orientativo.

    \item[\texttt{ExerciseModal.js}] Carga dinámicamente el catálogo de deportes
    desde el \textit{endpoint} \texttt{/api/sports} al abrirse, mostrando cada
    deporte junto con su ventana de peligro en horas. El formulario incluye también
    la duración de la sesión en minutos.
\end{description}

% ──────────────────────────────────────────────────────────────
\section{Diseño del perfil de usuario y parametrización clínica}
\label{sec:diseno-perfil}

La pantalla \texttt{ProfileScreen.js} permite al usuario configurar los parámetros
individuales de su tratamiento, que el agente de IA utilizará como contexto para
personalizar sus recomendaciones.

Los parámetros configurables son:

\begin{itemize}
    \item \textbf{Insulina rápida y lenta:} seleccionadas mediante \texttt{Picker}
          desplegable del catálogo de insulinas, mostrando nombre y duración (7 insulinas rápidas y 6 lentas, guardadas en base de datos. Las más comunes).
    \item \textbf{Ratios ICR por franja horaria:} un campo numérico por cada una de
          las cinco franjas del día (desayuno, media mañana, almuerzo, merienda y cena),
          expresado en gramos de carbohidratos por unidad de insulina (g/U).
    \item \textbf{Factor de Sensibilidad a la Insulina (FSI):} cuántos mg/dL reduce
          la glucosa una unidad de insulina rápida, utilizado para calcular el bolo
          de corrección cuando la glucosa está por encima del rango.
\end{itemize}

El guardado se realiza mediante una petición \texttt{PUT} a \texttt{/api/profile},
que ejecuta un \texttt{UPSERT} en la base de datos. El botón de guardar cambia
visualmente a verde con el texto ``Guardado correctamente'' durante 3 segundos tras una
operación exitosa, o a rojo con ``Error al guardar'' si la petición falla, antes
de volver a su estado neutro.

% ──────────────────────────────────────────────────────────────
\section{Diseño del agente de IA (lógica de recomendación según perfil)}
\label{sec:diseno-ia}

El agente de IA de GlucoMate es el componente más innovador del sistema. Su diseño
se basa en el paradigma de \textbf{\textit{context-aware prompting}}: en lugar de
enviar al modelo únicamente el mensaje del usuario, el backend construye en tiempo
real un \textit{system prompt} rico en datos clínicos del usuario antes de cada
llamada a la API.

\subsection{Recolección de contexto en tiempo real}
\label{subsec:contexto-ia}

Al recibir una petición \texttt{POST /api/chat}, el backend ejecuta cuatro consultas
SQL en paralelo mediante \texttt{Promise.all}:

\begin{enumerate}
    \item \textbf{Perfil de tratamiento:} ratios ICR por franja horaria y FSI.
    \item \textbf{Insulina activa (IOB):} suma de unidades de insulina rápida
          inyectadas en las últimas 4 horas.
    \item \textbf{Ejercicio en ventana de peligro:} deportes practicados cuya
          ventana de peligro aún no ha expirado, con su tipo de intensidad.
    \item \textbf{Glucosa actual:} último valor registrado y su tendencia.
\end{enumerate}

La ejecución en paralelo minimiza la latencia total de recopilación de contexto,
que de otro modo se acumularía de forma secuencial.

\subsection{Lógica de reducción por tipo de deporte}
\label{subsec:reduccion-deporte}

El sistema implementa una tabla de reducción de bolo en función del tipo de
intensidad del ejercicio practicado:

\begin{table}[H]
\centering
\caption{Porcentajes de reducción del bolo de insulina por tipo de intensidad
         deportiva, implementados en la constante \texttt{SPORT\_REDUCTION} del backend.}
\label{tab:sport-reduction}
\begin{tabular}{|l|c|p{6cm}|}
\hline
\textbf{Tipo de intensidad} & \textbf{Reducción} & \textbf{Comportamiento del agente} \\
\hline
Aeróbico        & 20\%  & Reduce el bolo calculado y lo explica. \\
Mixto           & 12,5\% & Reduce el bolo calculado y lo explica. \\
Baja intensidad & 5\%   & Reduce el bolo calculado y lo explica. \\
Anaeróbico      & 0\%   & No reduce el bolo, pero advierte del riesgo de hipoglucemia postejercicio. \\
\hline
\end{tabular}
\end{table}

Si el usuario ha practicado múltiples deportes en sus respectivas ventanas de peligro,
el sistema aplica la reducción máxima de entre todos ellos. Si alguno es anaeróbico,
se emite la advertencia correspondiente adicionalmente.

\subsection{Construcción del \textit{system prompt}}
\label{subsec:system-prompt}

El \textit{system prompt} que recibe \texttt{llama-3.3-70b-versatile} en cada llamada es un documento de texto
estructurado en secciones que incluye:

\begin{itemize}
    \item \textbf{Identidad y restricciones:} el agente se identifica como
          ``GlucoMate AI'', responde siempre en español y de forma concisa.
    \item \textbf{Datos en tiempo real:} glucosa actual, tendencia, ICR por franja
          horaria, FSI, y alerta de IOB activo si supera 1 UI.
    \item \textbf{Instrucciones de cálculo de bolo:} algoritmo paso a paso con
          fórmula explícita (\texttt{unidades = (g HC ÷ 10) × ICR}), con instrucción
          de solicitar la franja horaria si no está clara.
    \item \textbf{Reglas de glucosa actual:} si la glucosa es menor de 70, no
          recomendar insulina; si es mayor de 180, añadir bolo de corrección
          (\texttt{(glucosa - 100) ÷ FSI}); si está en rango, sin corrección.
    \item \textbf{Contexto de deporte:} la reducción de bolo calculada o la
          advertencia de deporte anaeróbico.
    \item \textbf{Aviso obligatorio:} el agente debe incluir al inicio de cada
          respuesta el texto ``Aviso: soy una IA de apoyo, no un médico. Esta sugerencia
          no reemplaza el criterio clínico.'', garantizando que el usuario siempre
          conoce las limitaciones del sistema.
\end{itemize}

El modelo se llama con \texttt{temperature: 0.2} para maximizar la consistencia
y el rigor del razonamiento clínico, minimizando la variabilidad creativa que
sería inapropiada en un contexto médico.

\subsection{Mecanismo de reintento}
\label{subsec:reintento-ia}

La llamada a la API de Groq está protegida por un mecanismo de reintento que
realiza hasta 3 intentos con una pausa de 5 segundos entre ellos ante errores 503
(\textit{Service Unavailable}) o 429 (\textit{Too Many Requests}). Cualquier otro
tipo de error se propaga inmediatamente al cliente.

% ──────────────────────────────────────────────────────────────
\section{Diseño de la interfaz de usuario}
\label{sec:diseno-ui}

El diseño visual de GlucoMate parte de los mockups digitales elaborados durante la
fase de prototipado y sigue los principios del \textit{Material Design} de Google,
que es la referencia de diseño nativa de Android.

Las decisiones de diseño más relevantes son:

\begin{itemize}
    \item \textbf{Jerarquía visual clara.} El valor de glucosa (110pt) domina
          visualmente la pantalla, ya que es la información más crítica para el
          usuario diabético. Los demás elementos (gráfica, botones de acción) se
          subordinan a él en tamaño e importancia visual.

    \item \textbf{Codificación de color semántica y consistente.} El sistema usa
          el mismo código de color en toda la interfaz: verde para comidas, rosa/rojo
          para insulina rápida, morado para insulina lenta, naranja para ejercicio.
          Esta consistencia permite al usuario interpretar la gráfica de forma
          intuitiva sin necesidad de leer etiquetas.

    \item \textbf{Botón NFC en posición de pulgar.} El botón de escaneo se sitúa
          en la esquina superior izquierda de la pantalla, zona de máxima
          accesibilidad con el pulgar derecho en dispositivos de pantalla grande.
          Cambia de color a naranja mientras espera el escaneo, proporcionando
          feedback visual inmediato del estado del sistema.

    \item \textbf{Navegación por pestañas inferior.} La barra de navegación
          inferior con dos pestañas (Glucosa y Mi Perfil) sigue el patrón de
          navegación estándar de Android y respeta el área segura del sistema
          mediante \texttt{useSafeAreaInsets}, evitando solapamientos con la
          barra de navegación del sistema.

    \item \textbf{Botón flotante de IA (FAB).} El acceso al chat de inteligencia
          artificial se implementa como un botón flotante de acción (FAB) de color
          azul marino (\texttt{\#1a237e}) con un icono de chispa (\texttt{sparkles}),
          que evoca la naturaleza inteligente de la funcionalidad sin ocupar espacio
          permanente en la interfaz principal.
\end{itemize}

% ──────────────────────────────────────────────────────────────
\section{Flujo global de interacción}
\label{sec:flujo-principal}

El flujo de uso más frecuente de GlucoMate —escanear el sensor y consultar al
agente de IA— sigue la siguiente secuencia:

\begin{figure}[H]
\centering
\resizebox{0.92\textwidth}{!}{%
\begin{tikzpicture}[
    lifeline/.style={font=\small, align=center},
    msg/.style={-{Latex[length=2.2mm]}, thick},
    reply/.style={msg, dashed},
    every node/.style={font=\scriptsize}
]
    \node[lifeline] (u) at (0,0) {Usuario};
    \node[lifeline] (f) at (3.4,0) {Frontend};
    \node[lifeline] (b) at (6.8,0) {Backend};
    \node[lifeline] (g) at (10.2,0) {Groq / BD};
    \foreach \x in {u,f,b,g} \draw[densely dashed] (\x.south) -- ++(0,-10.4);
    \draw[msg] (0,-0.8) -- node[above] {Pulsa NFC} (3.4,-0.8);
    \draw[msg] (3.4,-1.5) -- node[above] {NfcV} (6.8,-1.5);
    \draw[reply] (6.8,-2.2) -- node[above] {activo} (3.4,-2.2);
    \draw[msg] (3.4,-2.9) -- node[above] {GET Juggluco} (6.8,-2.9);
    \draw[reply] (6.8,-3.6) -- node[above] {lectura} (3.4,-3.6);
    \draw[msg] (3.4,-4.3) -- node[above] {POST /glucose} (6.8,-4.3);
    \draw[msg] (6.8,-5.0) -- node[above] {persistencia} (10.2,-5.0);
    \draw[reply] (10.2,-5.7) -- node[above] {ok + histórico} (6.8,-5.7);
    \draw[reply] (6.8,-6.4) -- node[above] {gráfica} (3.4,-6.4);
    \draw[msg] (0,-7.1) -- node[above] {consulta IA} (3.4,-7.1);
    \draw[msg] (3.4,-7.8) -- node[above] {POST /chat} (6.8,-7.8);
    \draw[msg] (6.8,-8.5) -- node[above] {prompt + contexto} (10.2,-8.5);
    \draw[reply] (10.2,-9.2) -- node[above] {respuesta} (6.8,-9.2);
    \draw[reply] (6.8,-9.9) -- node[above] {reply} (3.4,-9.9);
\end{tikzpicture}%
}
\caption{Diagrama de secuencia del flujo principal de GlucoMate: escaneo NFC y
         consulta al agente de IA.}
\label{fig:flujo-principal}
\end{figure}
